\begin{explanationbox}
	\textbf{\textcolor{CExplanation}{一句话直觉}}

	正弦定理的核心是：三角形中，每条边与其对角正弦的比值都相等，这个比值就是外接圆的直径，也就是 $2R$。在解三角形中，若已知条件足以确定三角形，通常可以结合内角和、正弦定理、余弦定理进行边角互化，求出其余边角；其中 SSA 情形可能出现无解、一解或两解，不能默认唯一。

	\medskip
	\textbf{\textcolor{CExplanation}{核心拆解}}
	\begin{description}[style=nextline, leftmargin=2.8em, labelsep=0.5em, itemsep=0.45em]
		\item[\textbf{要点一（边角比例关系）：}] 在同一个三角形中，边的长度和它所对角的正弦值成正比，即 $a \propto \sin A$，比例系数 $2R$ 对三条边都相同。
		\item[\textbf{要点二（外接圆桥梁）：}] 外接圆半径 $R$ 将三角形与圆联系起来，$2R$ 是直径。正弦定理的几何基础来自圆周角定理：一方面直径所对圆周角为直角，另一方面同弦所对圆周角相等或互补，从而保证相关角的正弦值相等。
	\end{description}

	\medskip
	\textbf{\textcolor{CExplanation}{几何本质（外接圆直径）}}

	过三角形一顶点作外接圆的直径，利用圆周角定理将角转移到直角三角形中，得到边长、外接圆直径与角的正弦之间的关系。因此，正弦定理是圆周角定理在边长计算上的代数表达。

	\medskip
	\textbf{\textcolor{CExplanation}{代数意义（边角互化工具）}}

	正弦定理给出了边 $a,b,c$ 与角 $A,B,C$ 以及外接圆半径 $R$ 之间的等式，从而可以在解三角形时进行边角互化。比如，把边的条件转化成角的关系，或把角的等式转化成边的比例，极大简化运算。但要注意，正弦定理在 SSA 情形下可能产生多解，不能只算出一个锐角就结束。

	\medskip
	\textbf{\textcolor{CExplanation}{考点价值（高频考法）}}
	\begin{itemize}[leftmargin=*, itemsep=0.5em, topsep=0.4em]
		\item \textbf{考法一（边角直接转换）：} 已知两角一边时，可利用正弦定理求其他边；已知两边及其中一边的对角（SSA）时，可先用正弦定理求另一角的正弦值，但必须进一步判断是否存在无解、一解或两解。
		\item \textbf{考法二（结合外接圆半径）：} 给出外接圆半径的条件，与正弦定理结合求边长或角。或者在证明题中通过边角互化推导三角形形状。
	\end{itemize}

	\medskip
	\textbf{\textcolor{CExplanation}{顿悟点}}

	\textbf{公式的灵魂是 $2R$}：只要出现 $a$ 和 $\sin A$，就想到连接它们的桥梁 $2R$。记清楚不是 $a/\sin A = R$，而是 $=2R$。使用边角转化时，注意不要丢掉系数 $2$。

	\medskip
	\textbf{\textcolor{CExplanation}{使用场景}}
	\begin{itemize}[leftmargin=*, itemsep=0.5em, topsep=0.4em]
		\item \textbf{场景一（解三角形基本量）：} 给出两个角和一个边（AAS，ASA），或两边和其中一边对角（SSA），求其他未知量。SSA 情形要特别检查多解。
		\item \textbf{场景二（边角互化证明）：} 在三角形中要证明边的关系或角的关系，常将边化为角的三角函数，或将角的三角函数化为边，再利用其他公式求解。
	\end{itemize}
\end{explanationbox}